% Options for packages loaded elsewhere
\PassOptionsToPackage{unicode}{hyperref}
\PassOptionsToPackage{hyphens}{url}
\documentclass[
  american,
]{article}
\usepackage{xcolor}
\usepackage{amsmath,amssymb}
\setcounter{secnumdepth}{-\maxdimen} % remove section numbering
\usepackage{iftex}
\ifPDFTeX
  \usepackage[T1]{fontenc}
  \usepackage[utf8]{inputenc}
  \usepackage{textcomp} % provide euro and other symbols
\else % if luatex or xetex
  \usepackage{unicode-math} % this also loads fontspec
  \defaultfontfeatures{Scale=MatchLowercase}
  \defaultfontfeatures[\rmfamily]{Ligatures=TeX,Scale=1}
\fi
\usepackage{lmodern}
\ifPDFTeX\else
  % xetex/luatex font selection
\fi
% Use upquote if available, for straight quotes in verbatim environments
\IfFileExists{upquote.sty}{\usepackage{upquote}}{}
\IfFileExists{microtype.sty}{% use microtype if available
  \usepackage[]{microtype}
  \UseMicrotypeSet[protrusion]{basicmath} % disable protrusion for tt fonts
}{}
\makeatletter
\@ifundefined{KOMAClassName}{% if non-KOMA class
  \IfFileExists{parskip.sty}{%
    \usepackage{parskip}
  }{% else
    \setlength{\parindent}{0pt}
    \setlength{\parskip}{6pt plus 2pt minus 1pt}}
}{% if KOMA class
  \KOMAoptions{parskip=half}}
\makeatother
\usepackage{color}
\usepackage{fancyvrb}
\newcommand{\VerbBar}{|}
\newcommand{\VERB}{\Verb[commandchars=\\\{\}]}
\DefineVerbatimEnvironment{Highlighting}{Verbatim}{commandchars=\\\{\}}
% Add ',fontsize=\small' for more characters per line
\newenvironment{Shaded}{}{}
\newcommand{\AlertTok}[1]{\textcolor[rgb]{1.00,0.00,0.00}{\textbf{#1}}}
\newcommand{\AnnotationTok}[1]{\textcolor[rgb]{0.38,0.63,0.69}{\textbf{\textit{#1}}}}
\newcommand{\AttributeTok}[1]{\textcolor[rgb]{0.49,0.56,0.16}{#1}}
\newcommand{\BaseNTok}[1]{\textcolor[rgb]{0.25,0.63,0.44}{#1}}
\newcommand{\BuiltInTok}[1]{\textcolor[rgb]{0.00,0.50,0.00}{#1}}
\newcommand{\CharTok}[1]{\textcolor[rgb]{0.25,0.44,0.63}{#1}}
\newcommand{\CommentTok}[1]{\textcolor[rgb]{0.38,0.63,0.69}{\textit{#1}}}
\newcommand{\CommentVarTok}[1]{\textcolor[rgb]{0.38,0.63,0.69}{\textbf{\textit{#1}}}}
\newcommand{\ConstantTok}[1]{\textcolor[rgb]{0.53,0.00,0.00}{#1}}
\newcommand{\ControlFlowTok}[1]{\textcolor[rgb]{0.00,0.44,0.13}{\textbf{#1}}}
\newcommand{\DataTypeTok}[1]{\textcolor[rgb]{0.56,0.13,0.00}{#1}}
\newcommand{\DecValTok}[1]{\textcolor[rgb]{0.25,0.63,0.44}{#1}}
\newcommand{\DocumentationTok}[1]{\textcolor[rgb]{0.73,0.13,0.13}{\textit{#1}}}
\newcommand{\ErrorTok}[1]{\textcolor[rgb]{1.00,0.00,0.00}{\textbf{#1}}}
\newcommand{\ExtensionTok}[1]{#1}
\newcommand{\FloatTok}[1]{\textcolor[rgb]{0.25,0.63,0.44}{#1}}
\newcommand{\FunctionTok}[1]{\textcolor[rgb]{0.02,0.16,0.49}{#1}}
\newcommand{\ImportTok}[1]{\textcolor[rgb]{0.00,0.50,0.00}{\textbf{#1}}}
\newcommand{\InformationTok}[1]{\textcolor[rgb]{0.38,0.63,0.69}{\textbf{\textit{#1}}}}
\newcommand{\KeywordTok}[1]{\textcolor[rgb]{0.00,0.44,0.13}{\textbf{#1}}}
\newcommand{\NormalTok}[1]{#1}
\newcommand{\OperatorTok}[1]{\textcolor[rgb]{0.40,0.40,0.40}{#1}}
\newcommand{\OtherTok}[1]{\textcolor[rgb]{0.00,0.44,0.13}{#1}}
\newcommand{\PreprocessorTok}[1]{\textcolor[rgb]{0.74,0.48,0.00}{#1}}
\newcommand{\RegionMarkerTok}[1]{#1}
\newcommand{\SpecialCharTok}[1]{\textcolor[rgb]{0.25,0.44,0.63}{#1}}
\newcommand{\SpecialStringTok}[1]{\textcolor[rgb]{0.73,0.40,0.53}{#1}}
\newcommand{\StringTok}[1]{\textcolor[rgb]{0.25,0.44,0.63}{#1}}
\newcommand{\VariableTok}[1]{\textcolor[rgb]{0.10,0.09,0.49}{#1}}
\newcommand{\VerbatimStringTok}[1]{\textcolor[rgb]{0.25,0.44,0.63}{#1}}
\newcommand{\WarningTok}[1]{\textcolor[rgb]{0.38,0.63,0.69}{\textbf{\textit{#1}}}}
\usepackage{graphicx}
\makeatletter
\newsavebox\pandoc@box
\newcommand*\pandocbounded[1]{% scales image to fit in text height/width
  \sbox\pandoc@box{#1}%
  \Gscale@div\@tempa{\textheight}{\dimexpr\ht\pandoc@box+\dp\pandoc@box\relax}%
  \Gscale@div\@tempb{\linewidth}{\wd\pandoc@box}%
  \ifdim\@tempb\p@<\@tempa\p@\let\@tempa\@tempb\fi% select the smaller of both
  \ifdim\@tempa\p@<\p@\scalebox{\@tempa}{\usebox\pandoc@box}%
  \else\usebox{\pandoc@box}%
  \fi%
}
% Set default figure placement to htbp
\def\fps@figure{htbp}
\makeatother
% definitions for citeproc citations
\NewDocumentCommand\citeproctext{}{}
\NewDocumentCommand\citeproc{mm}{%
  \begingroup\def\citeproctext{#2}\cite{#1}\endgroup}
\makeatletter
 % allow citations to break across lines
 \let\@cite@ofmt\@firstofone
 % avoid brackets around text for \cite:
 \def\@biblabel#1{}
 \def\@cite#1#2{{#1\if@tempswa , #2\fi}}
\makeatother
\newlength{\cslhangindent}
\setlength{\cslhangindent}{1.5em}
\newlength{\csllabelwidth}
\setlength{\csllabelwidth}{3em}
\newenvironment{CSLReferences}[2] % #1 hanging-indent, #2 entry-spacing
 {\begin{list}{}{%
  \setlength{\itemindent}{0pt}
  \setlength{\leftmargin}{0pt}
  \setlength{\parsep}{0pt}
  % turn on hanging indent if param 1 is 1
  \ifodd #1
   \setlength{\leftmargin}{\cslhangindent}
   \setlength{\itemindent}{-1\cslhangindent}
  \fi
  % set entry spacing
  \setlength{\itemsep}{#2\baselineskip}}}
 {\end{list}}
\usepackage{calc}
\newcommand{\CSLBlock}[1]{\hfill\break\parbox[t]{\linewidth}{\strut\ignorespaces#1\strut}}
\newcommand{\CSLLeftMargin}[1]{\parbox[t]{\csllabelwidth}{\strut#1\strut}}
\newcommand{\CSLRightInline}[1]{\parbox[t]{\linewidth - \csllabelwidth}{\strut#1\strut}}
\newcommand{\CSLIndent}[1]{\hspace{\cslhangindent}#1}
\ifLuaTeX
\usepackage[bidi=basic,shorthands=off]{babel}
\else
\usepackage[bidi=default,shorthands=off]{babel}
\fi
\ifLuaTeX
  \usepackage{selnolig} % disable illegal ligatures
\fi
\setlength{\emergencystretch}{3em} % prevent overfull lines
\providecommand{\tightlist}{%
  \setlength{\itemsep}{0pt}\setlength{\parskip}{0pt}}
\usepackage[letterpaper,margin=1in]{geometry}
\usepackage[scaled=0.95]{helvet}
\renewcommand{\familydefault}{\sfdefault}
\usepackage{xcolor}
\definecolor{ToolsBlue}{HTML}{2E74B5}
\definecolor{ToolsBlueDark}{HTML}{1F4D78}
\usepackage{sectsty}
\sectionfont{\color{ToolsBlue}}
\subsectionfont{\color{ToolsBlue}}
\subsubsectionfont{\color{ToolsBlueDark}}
\usepackage{microtype}
\usepackage{fancyhdr}
\pagestyle{fancy}
\fancyhf{}
\lhead{Tools Engineering Design Manual}
\rhead{Generated - Unapproved}
\cfoot{\thepage}
\usepackage{bookmark}
\IfFileExists{xurl.sty}{\usepackage{xurl}}{} % add URL line breaks if available
\urlstyle{same}
\hypersetup{
  pdftitle={Tools Engineering Design Manual},
  pdflang={en-US},
  hidelinks,
  pdfcreator={LaTeX via pandoc}}

\ifXeTeX
\special{pdf:trailerid [ <da2f2d9607739f8a6de4edbf14a13b98> <da2f2d9607739f8a6de4edbf14a13b98> ]}
\fi
\ifPDFTeX
\pdftrailerid{}
\pdftrailer{/ID [ <da2f2d9607739f8a6de4edbf14a13b98> <da2f2d9607739f8a6de4edbf14a13b98> ]}
\fi
\ifLuaTeX
\pdfvariable trailerid {[ <da2f2d9607739f8a6de4edbf14a13b98> <da2f2d9607739f8a6de4edbf14a13b98> ]}
\fi

\title{Tools Engineering Design Manual}
\usepackage{etoolbox}
\makeatletter
\providecommand{\subtitle}[1]{% add subtitle to \maketitle
  \apptocmd{\@title}{\par {\large #1 \par}}{}{}
}
\makeatother
\subtitle{Governed Shared Calculations and Interchange Pathways}
\author{}
\date{}

\begin{document}
\maketitle

\section{Release Status}\label{release-status}

This is the canonical editable QMD source for the Tools engineering
design manual. It is a governance scaffold, not a released engineering
manual. The file-level module inventory is complete for its declared
tracked-source denominator, but all detected calculation candidates
remain provisional and the stable calculation-pathway registry remains
empty. The required textbook chapter contract is qualified, but no
exemplar chapter is registered. Release remains blocked by TOOLS-D4
through TOOLS-D9. Generated HTML, LaTeX, PDF, and DOCX artifacts are
reproducible engineering records under TOOLS-D2, but none is approved
for publication or engineering use.

\section{Documentation Authority and Release
Boundary}\label{documentation-authority-and-release-boundary}

\subsection{Purpose}\label{purpose}

The manual will explain shared calculations and interchange contracts
from theory and assumptions through typed inputs, units, coordinate and
time conventions, algorithms, implementation symbols, tests,
uncertainty, limitations, consumer surfaces, and artifacts. Tools source
and evidence remain authoritative; the manual cannot approve a
calculation that its owning source, evidence, license, or human review
has not approved.

\subsection{Canonical Source}\label{canonical-source}

Only QMD under \texttt{manuals/tools} is editable manual source.
Existing user guides, ADRs, package documentation, generated API
references, and publications are separate governed products. They can
supply reviewed references, but none is a second mutable source for this
manual.

Generated LaTeX, PDF, DOCX, and HTML are non-editable artifacts. A
correction starts in canonical QMD, then regenerates every affected
representation through the qualified toolchain. TOOLS-D2 owns that
toolchain; TOOLS-D7 and TOOLS-D8 own artifact review and publication
evidence.

\subsection{Evidence and License
Classes}\label{evidence-and-license-classes}

Every calculation and claim must be classified as observed, derived,
model-scenario, external, provisional, unavailable, or approved within
an explicit scope. Every borrowed source and generated artifact must
preserve its license and attribution. Content from Tools\_Private or
another non-public source may not enter this public manual without
explicit public authority.

\subsection{Approval and Failure
Semantics}\label{approval-and-failure-semantics}

The default is deny-until-approved. Missing inventory coverage, source
authority, units, coordinate or time convention, implementation symbol,
test, validity range, uncertainty, license, artifact digest, page
review, accessibility review, or human approval keeps the affected item
blocked. Unsupported values are never replaced with synthetic or
plausible values without an explicit model-scenario contract.

\subsection{Exemptions}\label{exemptions}

Freshness failures are release-blocking. The TOOLS-D1 module baseline is
recomputed from tracked implementation and governed configuration files
by
\texttt{python\ -m\ scripts.build\_tools\_module\_inventory\ -\/-check};
later TOOLS-D6 work must extend freshness from the file baseline to
stable calculation pathways. Future exemptions must be structured,
narrowly scoped, justified, owned, reviewed, and expiring. An exemption
cannot authorize an equation, scientific claim, artifact, operating use,
license transfer, or public projection.

\subsection{Current Boundary}\label{current-boundary}

TOOLS-D1 establishes a complete file-level denominator for its declared
tracked scope. TOOLS-D2 establishes deterministic byte reproduction and
governed semantic preservation across HTML, LaTeX, PDF, and DOCX. Its
generated artifacts remain non-editable and unapproved. The conservative
inventory does not establish stable calculation identity, equation
correctness, validity ranges, uncertainty, page/accessibility approval,
or publication authority. Those pathways remain empty and blocked
pending TOOLS-D4 through TOOLS-D9.

\section{Module and Calculation-Candidate
Inventory}\label{module-and-calculation-candidate-inventory}

\subsection{Purpose and Denominator}\label{purpose-and-denominator}

\texttt{manifests/module-inventory.json} is the deterministic TOOLS-D1
index for every tracked implementation source file and governed
configuration module in the repository. It scans recognized Python,
Python stub, JavaScript, TypeScript, Rust, MATLAB, C/C++, Arduino,
shell, PowerShell, HTML, CSS, Prisma, JSON, YAML, TOML, INI, XML, URDF,
and web-manifest files. Tests are evidence inputs rather than
implementation rows. Generated, vendored, archived, experimental,
example, documentation, resource, and manual trees are excluded
explicitly in the manifest, so a consumer can reproduce the exact
denominator. The index binds bounded entry shards by normalized SHA-256,
contiguous index, entry count, and first/last path. Both the index and
every shard remain below the repository's one-megabyte file budget;
consumers reassemble one logical view only after verifying every shard
and the exact declared shard set.

At this baseline, 3,437 modules are classified: 808 conservative
calculation candidates and 2,629 non-calculation modules. One legacy
JavaScript file has a non-UTF-8 encoding and remains explicitly blocked
rather than being decoded by an implicit platform default. These counts
describe inventory coverage only; they are not test coverage,
calculation approval, or publication readiness.

\subsection{Deterministic Identity and
Integrity}\label{deterministic-identity-and-integrity}

Each file receives a provisional path-derived identifier and a SHA-256
digest of normalized text bytes. The normalization maps CRLF and bare CR
to LF before hashing. For module path (p\_i) and normalized content
digest (h\_i), the tree authority is

\[
H_{\mathrm{tree}} = \mathrm{SHA256}\left(
  \Vert_{i=1}^{n} p_i : h_i
\right),
\]

where rows are ordered by normalized repository-relative path. This
contract makes the result independent of Windows checkout line endings
while retaining content changes. The path-derived identifier is not a
stable calculation ID; TOOLS-D3 owns the structural chapter contract;
TOOLS-D4 owns exemplar registration and stable equation-to-code
identities that survive file movement.

\subsection{Classification Contract}\label{classification-contract}

A module is a \texttt{calculation} candidate only when its normalized
path contains a governed calculation marker, such as dynamics,
kinematics, pressure drop, signal processing, geometry, solver, energy,
force, or trajectory, or when the source imports a recognized scientific
mathematics library. All other modules are \texttt{non-calculation}.
This intentionally conservative static screen can have false positives
and false negatives, so every calculation candidate has
\texttt{authority\_status=provisional} and
\texttt{review\_status=review-required}.

The inventory also records language, purpose where statically available,
maintainer family, public surfaces and routes, mapped tests and
validation, ADRs, equation references, literal units, citations,
artifact digests, risk tags, normalized byte count, and content digest.
Missing tests, sources, units, chapters, equation pathways, validation,
routes, or artifact digests remain \texttt{unavailable},
\texttt{not-applicable}, or explicitly pending their owning subepic;
they are never synthesized. Non-calculation status means only that the
conservative detector found no calculation signal. It does not waive
public-API, security, licensing, accessibility, or downstream
compatibility review.

\subsection{Consumer and Failure
Contract}\label{consumer-and-failure-contract}

The repository-owned schema is
\texttt{schemas/module-inventory.schema.json}, version
\texttt{tools-module-inventory/1.0.0}. It extends the program-owned
calculation registry without copying or competing with that shared
schema. Both JSON Schema validation and the Python consumer reject
unknown fields, unsupported versions, unsafe paths, invalid hashes,
duplicate identifiers, unsorted duplicate traceability arrays,
denominator mismatches, and unsupported state values.

Any tracked in-scope module addition, removal, rename, or normalized
content change makes inventory check mode fail until the manifest is
regenerated and reviewed. A green inventory check proves only
deterministic classification and freshness of this file-level baseline.
Public projection remains blocked until the remaining stable-pathway,
page, accessibility, licensing, publication, and human-approval gates
are complete.

\section{Reproducible Rendering
Contract}\label{reproducible-rendering-contract}

\subsection{Outcome and Authority}\label{outcome-and-authority}

TOOLS-D2 renders the canonical QMD into deterministic HTML, LaTeX, PDF,
and DOCX representations. Generated outputs are engineering records, not
approval. Only QMD may be corrected directly. Every output remains
blocked from public projection until semantic, page, accessibility,
provenance, licensing, and human-review gates owned by later subepics
are complete.

The renderer uses the pinned command versions and inputs recorded in
\texttt{toolchain-lock.json}. It verifies the exact first version line
before render, expands only bounded QMD includes, uses the governed
bibliography and reference DOCX, disables TeX shell escape, normalizes
volatile package metadata, and binds every representation into
\texttt{manifests/artifacts.json}. The command model follows the
documented Pandoc conversion contract (Pandoc contributors 2026) while
QMD remains compatible with the pinned Quarto authoring surface (Posit
Software, PBC 2026).

\begin{figure}
\centering
\includegraphics[width=0.8\linewidth,height=\textheight,keepaspectratio,alt={Deterministic rendering pathway.}]{figures/render-pipeline.png}
\caption{Deterministic rendering pathway.}
\end{figure}

\subsection{Semantic, Unit, and Warning
Preservation}\label{semantic-unit-and-warning-preservation}

The semantic contract is deliberately format-neutral. Every
representation must retain the governed title, release boundary,
representative unit string, warning, figure label, and bibliography
identity before receiving a shared semantic digest. Example unit string:
10 m/s.

WARNING: Provisional calculations remain blocked.

This gate proves preservation of required meaning across the four
representations. It does not prove calculation correctness, equation
coverage, accessibility, page aesthetics, or release approval. Missing
text, a stale digest, an unpinned tool, an unsafe include, or an absent
artifact fails closed.

\subsection{Reproduction and Consumer
Contract}\label{reproduction-and-consumer-contract}

Run
\texttt{python\ -m\ scripts.render\_tools\_design\_manual\ -\/-verify-toolchain}
before a qualified build. Run
\texttt{python\ -m\ scripts.render\_tools\_design\_manual\ -\/-render}
to replace renderer-owned artifacts, then
\texttt{python\ -m\ scripts.render\_tools\_design\_manual\ -\/-check} to
validate source, toolchain, byte, and semantic bindings. The strict JSON
Schemas allow consumers to validate the lock and manifest without
importing producer code.

Artifact rows expose media type, byte count, byte digest, and shared
semantic digest. The manifest also preserves the owning epic, review
owner, explicit blockers, and a null immutable source commit while this
local stack is not a merged authority. A consumer must reject unknown
fields, versions, formats, paths, hashes, status promotion, or semantic
divergence.

\section{References}\label{references}

\section{Required Textbook Chapter
Contract}\label{required-textbook-chapter-contract}

\subsection{Outcome and Authority}\label{outcome-and-authority-1}

TOOLS-D3 defines the public structure that every calculation-level
textbook chapter must satisfy before registration. The contract is
structural and traceability-focused: passing it does not approve an
equation, implementation, dataset, artifact, license, or engineering
use. The registry is intentionally empty until TOOLS-D4 supplies
reviewed exemplar chapters.

Two version identifiers are fixed:

\begin{itemize}
\tightlist
\item
  Contract: \texttt{tools-textbook-chapter-contract/1.0.0}
\item
  Registry: \texttt{tools-textbook-chapter-registry/1.0.0}
\end{itemize}

Their repository-owned JSON Schemas allow external consumers to validate
shape without importing producer code. The Python consumer then applies
stricter invariants for normalized paths, ordering, uniqueness,
authority state, and QMD content.

\subsection{Required Chapter Anatomy}\label{required-chapter-anatomy}

Fourteen ordered sections are mandatory. A registered chapter must
present, in order: Purpose and Scope; Design by Contract; Coordinates,
Frames, and Time; Units and Dimensions; Derivation and Assumptions;
Algorithm; Implementation Symbols; Error and Failure Semantics;
Uncertainty and Sensitivity; Verification and Validation; Limits and
Applicability; Worked Examples; References and Provenance; and Revision
History.

The contract also declares mandatory third-level headings. For example,
Design by Contract requires Preconditions, Postconditions, and
Invariants. Coordinates requires frame, handedness, and time
declarations. Verification and Validation requires test evidence and
acceptance criteria. Exact ordered headings remove format-dependent
inference and make an omitted obligation a deterministic failure.

\subsection{Registration and
Traceability}\label{registration-and-traceability}

Each registered chapter has one stable chapter ID and one normalized QMD
path. Its descriptor binds the chapter version and provisional authority
to an owner, review owner, stable calculation IDs, module IDs,
verification tests, source citations, explicit limitations, and revision
history. All traceability arrays are non-empty, unique, and sorted;
duplicate chapter IDs or paths are rejected.

An \texttt{approved} chapter claim is unsupported by this D3 contract
because later subepics own page, accessibility, publication, and
human-review evidence. D3 accepts only \texttt{provisional} and
\texttt{verified-unapproved} chapter descriptors. The manual and every
generated representation therefore remain \texttt{generated-unapproved}.

\subsection{Linter Algorithm and
Evidence}\label{linter-algorithm-and-evidence}

The linter loads the contract and registry through strict typed
consumers, resolves each registered path beneath
\texttt{manuals/tools/chapters}, decodes UTF-8, and rejects forbidden
placeholder or private-source tokens. It extracts level-2 and level-3
headings, compares them with the exact ordered contract, and checks that
every required section has governed body content.

For normalized chapter text \(T\) and ordered section blocks \(S_i\),
the linter records deterministic evidence as

\[
h_T = \operatorname{SHA256}(\operatorname{LF}(T)), \qquad
h_i = \operatorname{SHA256}(\operatorname{LF}(S_i)).
\]

These hashes detect byte-level content drift while remaining independent
of Windows or POSIX checkout line endings. They are lint evidence, not
scientific approval or an immutable publication digest.

\subsection{Failure and Ownership
Boundaries}\label{failure-and-ownership-boundaries}

Missing files, invalid UTF-8, unsafe paths, unknown fields or versions,
duplicate identities, unsorted traceability, missing or reordered
headings, empty governed content, placeholders, private-source
references, unsupported status promotion, and stale generated artifacts
all fail closed. The linter does not synthesize equations, units,
sources, examples, tests, or limitations.

TOOLS-D4 owns exemplar chapter registration. TOOLS-D5 expands phased
module coverage, TOOLS-D6 owns executable freshness, TOOLS-D7 owns
rendered and accessibility review, and TOOLS-D8 owns immutable public
projection and human approval.

\subsection{Consumer Workflow}\label{consumer-workflow}

Authors edit only canonical QMD and the machine-readable chapter
registry. Run these commands in order:

\begin{Shaded}
\begin{Highlighting}[]
\NormalTok{python {-}m scripts.lint\_tools\_textbook\_chapters}
\NormalTok{python {-}m scripts.render\_tools\_design\_manual {-}{-}render}
\end{Highlighting}
\end{Shaded}

They regenerate and check all four representations. The combined
governance, inventory, chapter, and artifact freshness gates must pass
before review.

\protect\phantomsection\label{refs}
\begin{CSLReferences}{1}{1}
\bibitem[\citeproctext]{ref-pandoc-manual}
Pandoc contributors. 2026. \emph{Pandoc User's Guide}. Pandoc.
\url{https://pandoc.org/MANUAL.html}.

\bibitem[\citeproctext]{ref-quarto-guide}
Posit Software, PBC. 2026. \emph{Quarto Guide}. Posit Software, PBC.
\url{https://quarto.org/docs/guide/}.

\end{CSLReferences}

\end{document}
